\documentclass[quiet]{ctexart}

\ExplSyntaxOn \makeatletter

\keys_define:nn { hduthesis / exam }
  {
    watermark .tl_set:N = \l__hduthesis_exam_watermark_tl,
  }
\NewDocumentCommand \hdutexset {m} { \keys_set:nn { hduthesis / exam } {#1} }
\DeclareHookRule { shipout / background }
  { hduthesis / exam } { before } { pgfrcs }
\AddToHook { shipout / background } [ hduthesis / exam ]
  {
    \tl_if_empty:NF \l__hduthesis_exam_watermark_tl
      {
        \tikz [ remember~picture, overlay ]
          \node [ opacity = .2 ] at (current~page)
            {
              \includegraphics [ width = \linewidth ]
                { \l__hduthesis_exam_watermark_tl }
            };
      }
  }

\RequirePackage
  {
    geometry, mathtools, array, bm, esint, keytheorems, libertine, inconsolata,
    siunitx, fixdif, derivative, physics2, tasks, anyfontsize
  }
\RequirePackage[inline]{enumitem}
\RequirePackage{fontawesome5, fancyqr, circuitikz}
\hook_gput_code:nnn { cmd/faIcon/before } { . }
  {
    \group_begin:
    \fontsize { \fp_eval:n { .9 * \f@size } } { 0 } \selectfont
  }
\hook_gput_code:nnn { cmd/faIcon/after } { . } { \group_end: }
\geometry{margin = 1in}

\RequirePackage[svgnames]{xcolor}
\newkeytheorem{note}
  [ style = remark, numbered = false,
    name = \color{DarkRed}\sffamily \faIcon{info-circle}~Tip. ]
\newkeytheorem{problem}
  [ style = definition,
    name = \color{DarkGreen}\sffamily \faIcon{pen-square}~Problem ]
\newkeytheorem{solution}
  [ style = definition, numbered = false,
    name = \color{MidnightBlue}\sffamily \faIcon{check-square}~Solution ]
\usephysicsmodule{ab, op.legacy}
\usephysicsmodule{bm-um.legacy}
\settasks { label = \sffamily\Alph*., item-indent = 2.5em, label-offset = .5pt }
\newcommand \true {\addtocounter{task}{1}\faIcon{check-circle}}
\def \@result#1 {\underline{~#1~}}
\def \s@result {\textsf{【\qquad 】}}
\DeclareRobustCommand \result {\@ifstar\s@result\@result}
\keys_define:nn { textfig }
  {
    ratio      .tl_set:N    = \l__textfig_ratio_tl,
      ratio    .initial:n   = 7 \c_colon_str 2,
    position   .tl_set:N    = \l__textfig_position_tl,
      position .initial:n   = c,
    unknown    .code:n      = \tl_if_novalue:nF {#1}
      { \tl_set_eq:NN \l__textfig_position_tl \l_keys_key_tl }
  }
\NewDocumentEnvironment{textfig}{ O{} }
  {
    \group_begin:
    \keys_set:nn { textfig } {#1}
    \seq_set_split:Nee \l__whats_textfig_ratio_seq { \c_colon_str }
      { \l__textfig_ratio_tl }
    \noindent
    \minipage [ \l__textfig_position_tl ]
      {
        \fp_eval:n
          {
            .96*\seq_item:Nn \l__whats_textfig_ratio_seq {1}/
              (
                \seq_item:Nn \l__whats_textfig_ratio_seq {1} + 
                \seq_item:Nn \l__whats_textfig_ratio_seq {2}
              )
          }\dimexpr \linewidth
      }
  }
  {
    \endminipage
    \vspace{4.5pt plus 1.5pt minus 1.5pt} \par
    \group_end:
  }
\NewDocumentCommand \textfigsplit {}
  {
    \endminipage
    \hspace{.02\linewidth}
    \minipage [ \l__textfig_position_tl ]
      {
        \fp_eval:n
          {
            .96*\seq_item:Nn \l__whats_textfig_ratio_seq {2}/
              (
                \seq_item:Nn \l__whats_textfig_ratio_seq {1} + 
                \seq_item:Nn \l__whats_textfig_ratio_seq {2}
              )
          }\dimexpr \linewidth
      }
  }

% \hook_gput_code:nnn { cmd/maketitle/before } { . } { \vskip -\baselineskip }
% \hook_gput_code:nnn { cmd/maketitle/after } { . }
%   { \__weyoung_maketitle_aftercode_aux: }
\cs_new_protected:Nn \__weyoung_maketitle_aftercode_aux:
  {
    \vspace*{-2\baselineskip}
    \begin{center}
      \begin{minipage}{.66\textwidth}
        \centering
        \begin{tasks}
          [ label-width = 5\ccwd, item-indent = 2cm,
            after-item-skip = 0pt ](2)
          \task[\textbf{课程信息}：] 理学院/A0715022
          \task[\textbf{授课教师}：] 陈许敏
          \task[\textbf{解析提供}：] \textsf{Skyrmion}
          \task[{\makebox[5em][s]{\textbf{\textsf\TeX~排版}：}}] \textsf{Skyrmion}
        \end{tasks}
      \end{minipage}
      \begin{minipage}{.15\textwidth}
        \centering
        \begin{tikzpicture}
          \node [ inner~sep = 0 ] (wechat) {
            \fancyqr
              [ size = 4em, image~padding = 1,
                image = \raisebox{4em/33}{\textcolor{DarkGreen}{\faIcon{weixin}}},
                r~color = DarkGreen, version = 4
              ]{http://weixin.qq.com/r/sRJiemTEGjllrSAZ90ej}
            };
          \node [ inner~sep = .5ex, align = center, right, font = \scriptsize ]
            at (wechat.east) {未\\央\\学\\社};
        \end{tikzpicture}
      \end{minipage}
      \begin{minipage}{.15\textwidth}
        \centering
        \begin{tikzpicture}
          \node [ inner~sep = 0 ] (qq) {
            \fancyqr
              [ size = 4em, image~x~padding = 2, image~y~padding = 1,
                image = \raisebox{4em/33}{\textcolor{MidnightBlue}{\faIcon{qq}}},
                l~color = MidnightBlue, version = 4
              ]{https://qm.qq.com/q/PH7CsQWvm2}
            };
          \node [ inner~sep = .5ex, align = center, left, font = \scriptsize ]
            at (qq.west) {物\\理\\营};
        \end{tikzpicture}
      \end{minipage}
    \end{center}
  }

\newcommand \iu {{\mathup i\mkern1mu}}
\newcommand \eu {\mathup e}
\AtBeginDocument{%
  \let \Re \relax \let \Im \relax
  \DeclareMathOperator{\Re}{Re}
  \DeclareMathOperator{\Im}{Im}
}

\makeatother \ExplSyntaxOff
\usetikzlibrary{ fadings, patterns, arrows, decorations.markings }
\tikzset
  {
    > = stealth',
    every picture/.append style =
      { line join = round, line cap = round, thick }
  }

\setCJKmainfont{LXGW WenKai}[AutoFakeSlant]
\setCJKsansfont{LXGW Marker Gothic}[AutoFakeBold, AutoFakeSlant, Scale = .95]
\setCJKmonofont{LXGW WenKai Mono}[AutoFakeBold, AutoFakeSlant]
\ctexset{section/format+ = \raggedright}
\usepackage[warnings-off={mathtools-colon,mathtools-overbracket}]{unicode-math}
\setmathfont[AutoFakeBold]{STIX Two Math}

\makeatletter \def \@result #1{\underline{ \quad }} \makeatother
\DeclareMathOperator{\DFT}{DFT}

\title{\sffamily\bfseries 24Spring \textsc{Hdu} 「数理方法」期末试题}
\author{} \date{}

\begin{document}

\maketitle

\section{选择题（每题 4 分，共 16 分）}

\begin{problem}
  计算下列表达式：
  $\sin\varphi + \sin2\varphi + \sin3\varphi + \cdots + \sin n\varphi$
  （$\varphi$ 为实函数），\result*项正确.
  \begin{tasks}(2)
    \task $\Im\frac{\exp(\iu\varphi) - \exp[\iu(n + 1)\varphi]}{1 - \exp(\iu\varphi)}$
    \task $\frac{\sin(\varphi/2) - \sin(n + 1/2)\varphi}{2\cos(\varphi/2)}$
    \task $\frac{\cos(\varphi/2) - \cos(n + 1/2)\varphi}{2\sin(\varphi/2)}$
    \task $\Re\frac{\exp(\iu\varphi) - \exp[\iu(n + 1)\varphi]}{1 - \exp(\iu\varphi)}$
  \end{tasks}
\end{problem}
\begin{solution}

\end{solution}

\begin{problem}
  已知解析函数 $f(z)$ 的实部 $u = x^3 - 6x^2y - 3xy^2 -2y^3$, $f(0) = 0$. 该解析函数为
  \begin{tasks}(4)
    \task $z^3$
    \task $z^3(1 - 2\iu)$
    \task $z^3(1 + 2\iu)$
    \task $2\iu z^3$
  \end{tasks}
\end{problem}
\begin{solution}

\end{solution}

\begin{problem}
  在圆柱坐标中，一点的位置 $(4,2\pi/3,3)$ 给出，则该点在球坐标系中的坐标为
  \begin{tasks}(4)
    \task $(4,2\pi/3,3)$
    \task $(0,-4,3)$
    \task $(5,\arcsin(4/5),2\pi/3)$
    \task 以上都不对
  \end{tasks}
\end{problem}
\begin{solution}

\end{solution}

\begin{problem}
  矢量场的散度满足的关系为
  \begin{tasks}(2)
    \task $\nabla \cdot \bm A = \lim\limits_{\Delta \nu \to 0} 
           \frac{\oiint \bm A \cdot \d\bm s}{\Delta \nu}$
    \task $\nabla \cdot \bm A = \pdv{A_x}x + \pdv{A_y}y + \pdv{A_z}z$
    \task $\nabla \cdot \bm A \d \nu = \oiint_s \bm A \cdot \d \bm s$
    \task 以上都正确
  \end{tasks}
\end{problem}
\begin{solution}
  
\end{solution}

\section{填空题（每题4分，共 8 分）}

\begin{problem}
  计算下列数值（$x$ 为常数）：$\cos\iu x = $\result{}.
\end{problem}
\begin{solution}
  
\end{solution}

\begin{problem}
  函数 $f(z) = \frac{1}{1 + z^{3\pi}}$ 的奇点及其对应的留数是\result{}.
\end{problem}
\begin{solution}
  
\end{solution}

\section{解答题}

% 3
\begin{problem}[本题 8 分]
  设 $f(z) = my^3 + nx^2y - \iu(x^3 - 3xy^2)$ 为解析函数，试确定 $m$, $n$, $l$.
\end{problem}
\begin{solution}
  
\end{solution}

% 4
\begin{problem}[本题 10 分]
  分别沿路径\begin{enumerate*}
    \item $(-1,0)$ 到 $(1,0)$ 的直线段，路径
    \item 单位圆周的上半部分，路径
    \item 单位圆周的下半部分，
  \end{enumerate*}
  计算积分 $\int_{-1}^1|z| \d z$ 的值.
\end{problem}
\begin{solution}
  
\end{solution}

% 5
\begin{problem}[本题 4 分]
  写出解析函数在任一点处函数值用境界线用原积分表示的科西公式，以及解析函数的二阶导数的科西公式.
\end{problem}
\begin{solution}
  
\end{solution}

% 6
\begin{problem}[本题 11 分]\leavevmode
  \begin{enumerate}
    \item 在 $z_0 = \iu$ 的邻域上把 $\sqrt[3]z$ 展开为泰勒级数.
    \item 将函数 $f(z) = \frac{z^2 - 2z + 5}{(z^2 + 1)(z - 2)}$ 在下列圆环域内展开成洛朗级数：$1 < |z| < 2$.
  \end{enumerate}
\end{problem}
\begin{solution}
  
\end{solution}

% 7
\begin{problem}[本题 6 分]
  已知矢量 $\bm E = \bm e_x(3y - az) + \bm e_y(bx - 2z) - \bm e_z(cy + z)$，当 $E$ 为无源场时，求 $a$, $b$, $c$ 的值，
  并计算通量积分 $\oint_S\bm E \cdot \d\bm S$，其中 $\bm S$ 是顶点在原点的单位立方体表面.
\end{problem}
\begin{solution}
  
\end{solution}

% 8
\begin{problem}[本题 10 分]
  应用留数定理计算下列实变函数的定积分问题
  \begin{enumerate}
    \item $\int_0^{2\pi} \frac{\d x}{(1 + \epsilon \cos x)^2}$ ($0 < \epsilon < 1$)
    \item $\int_0^{+\infty} \frac{1 + x^2}{1 + x^6} \d x$.
  \end{enumerate}
\end{problem}
\begin{solution}
  
\end{solution}

% 9
\begin{problem}[本题 5 分]
  确定函数 $f(z) = \frac{z\eu^z}{(z - a)^3}$ 的奇点，并求出函数在奇点的留数.
\end{problem}
\begin{solution}
  
\end{solution}

% 10
\begin{problem}[本题 8 分]
  利用达朗贝尔公式求解无限长弦的自由振动，设弦的初始位移为 $\varphi(x)$，初始速度为 $-a\varphi(x)$.
\end{problem}
\begin{solution}
  
\end{solution}

% 11
\begin{problem}[本题 4 分]
  把幂级数 $\sum_{k = 0}^\infty a_k (z - z_0)^k = a_0 + a_1(z - z_0) + a_2(z - z_0)^2 + \ldots$ 逐项积分，求所得级数的收敛半径，以此验证逐项积分并不改变收敛半径.
\end{problem}
\begin{solution}
  
\end{solution}

% 12
\begin{problem}[本题 10 分]
  利用分离变数法求解细杆导热问题 ($u_t - a^2u_{xx} = 0$). 杆长 $L$，两端保持为零度，初始温度分布为
  \[u\big|_{t = 0 }= bx(l - x)/l^2\]
\end{problem}
\begin{solution}
  
\end{solution}

\end{document}